\documentclass[a4paper]{article}

%------------------------------------------------
% الحزم
%------------------------------------------------

\usepackage{enumitem}
\usepackage[margin=2.5cm]{geometry}
\usepackage{amsmath,amsfonts,amssymb}
\usepackage[charter]{mathdesign}
\usepackage{graphicx}
\usepackage{array}
\usepackage{xcolor}
\usepackage{needspace}
\usepackage{fancyhdr}
\usepackage{polyglossia}

%------------------------------------------------
% اللغات والخطوط
%------------------------------------------------

\setdefaultlanguage[
    calendar=gregorian,
    numerals=maghrib
]{arabic}

\setotherlanguage{english}

% إذا لم تكن الخطوط مثبتة في الجهاز، تُؤخذ من المجلد ../../fonts/
\IfFontExistsTF{Amiri}{%
    \newfontfamily\arabicfont[Script=Arabic, Scale=1.2]{Amiri}
}{%
    \newfontfamily\arabicfont[
        Script=Arabic, Scale=1.2,
        Path=../../fonts/, Extension=.ttf,
        UprightFont=*-Regular, BoldFont=*-Bold,
        ItalicFont=*-Italic, BoldItalicFont=*-BoldItalic
    ]{Amiri}
}

\IfFontExistsTF{AlBattar}{%
    \newfontfamily\arabicfontsf[Script=Arabic, Scale=1.3]{AlBattar}
    \newfontfamily\khatt[Script=Arabic, Scale=1.3]{AlBattar}
}{%
    \newfontfamily\arabicfontsf[
        Script=Arabic, Scale=1.3, Path=../../fonts/,
        BoldFont=ae_AlBattar.ttf
    ]{ae_AlBattar.ttf}
    \newfontfamily\khatt[
        Script=Arabic, Scale=1.3, Path=../../fonts/,
        BoldFont=ae_AlBattar.ttf
    ]{ae_AlBattar.ttf}
}

%------------------------------------------------
% نسخة الطالب / نسخة التصحيح
% لإنتاج نسخة التصحيح: عرّف \withanswers قبل \input لهذا الملف
%------------------------------------------------

\newif\ifcorrection
\ifdefined\withanswers \correctiontrue \else \correctionfalse \fi

\newcommand{\emptybox}{$\square$}
\ifcorrection
    \newcommand{\correctbox}{\textcolor{green!50!black}{$\boxtimes$}}
    \newcommand{\correct}[1]{\textcolor{green!40!black}{\textbf{#1}}}
    \newcommand{\explanation}[1]{%
        \par\vspace{3pt}
        \fcolorbox{gray!60}{gray!8}{\parbox{\dimexpr\linewidth-2\fboxsep-2\fboxrule}{%
        \small\textbf{التعليل:} #1}}}
\else
    \newcommand{\correctbox}{$\square$}
    \newcommand{\correct}[1]{#1}
    \newcommand{\explanation}[1]{}
\fi

\newcommand{\pts}[1]{\textbf{(#1 ن)}}

%------------------------------------------------
% إعداد الصفحة
%------------------------------------------------

\renewcommand{\footrulewidth}{1pt}
\renewcommand{\headrulewidth}{1pt}

\rhead{\khatt المدرسة العليا للأساتذة ببوزريعة}
\lhead{\khatt ملحقة بومرداس}

\cfoot{\thepage}
\lfoot{\khatt الأستاذ: د. بن عبد الرزاق يوسف}

\parindent 0pt
\setlength{\headsep}{15pt}
\arraycolsep=1.4pt

\pagestyle{fancy}

\setlist[enumerate,1]{label=\textbf{السؤال \arabic*:}, leftmargin=*, itemsep=12pt, align=left}

%================================================
% بداية الوثيقة
%================================================

\begin{document}

\begin{center}

{\Large\khatt\textbf{الواجب الثالث}\ifcorrection\ \textbf{-- التصحيح النموذجي}\fi}

\vspace{0.2cm}

{\Large\khatt\textbf{في الذكاء الاصطناعي}}

\end{center}

\rule{\textwidth}{1pt}

\ifcorrection\else
\textbf{الاسم واللقب:} \dotfill \quad \textbf{الفوج:} \makebox[3cm]{\dotfill}

\vspace{0.3cm}
\fi

\begin{center}

{\khatt\textbf{أسئلة متعددة الاختيارات (20 سؤالا -- 15 نقطة)}}

\end{center}

\rule{\textwidth}{1pt}

\textbf{التعليمة:} ضع علامة \textbf{$\times$} في المربع المقابل للإجابة الصحيحة. لكل سؤال إجابة صحيحة واحدة فقط.

\vspace{0.3cm}

\begin{enumerate}

\needspace{9\baselineskip}\item
\textbf{كتب طالب: «أعطني خطة للمراجعة»، وكتب زميله: «تصرّف كمساعد أكاديمي. أنشئ خطة مراجعة لمادة النحو مدتها 7 أيام، بمعدل ساعتين يوميًا. حدّد موضوع كل يوم، وأضف تمرينًا قصيرًا، ونظّم الخطة في جدول». لماذا يُتوقع أن يحصل الثاني على إجابة أفضل؟} \hfill \pts{0,75}
\begin{itemize}[leftmargin=1.6cm, labelsep=0.3cm, itemsep=2pt, topsep=4pt]
    \item[\correctbox\ \textbf{أ)}] \correct{لأنه أعطى الأداة كل ما تحتاجه لتفهم طلبه: دورًا ومهمة وسياقًا وقيودًا وتنسيقًا واضحًا للإجابة.}
    \item[\emptybox\ \textbf{ب)}] لأن الأمر الأطول يحصل دائمًا على إجابة أفضل، مهما كانت المعلومات التي يحتويها أو طريقة ترتيبها.
    \item[\emptybox\ \textbf{ج)}] لأن عبارة «تصرّف كمساعد أكاديمي» كلمة سرّ خاصة تفتح خصائص متقدمة لا تظهر لبقية المستخدمين.
    \item[\emptybox\ \textbf{د)}] لأن الأداة تعرف مسبقًا مستوى كل طالب ومادته، فتكافئ من يكتب بلغة رسمية بإجابات أكثر تفصيلًا.
\end{itemize}
\explanation{الأداة تعرف فقط ما نكتبه لها. الأمر الثاني يحدد الدور والمهمة والسياق (النحو) والقيود (7 أيام، ساعتان) والتنسيق (جدول)، فتفهم الأداة المطلوب بدقة.}

%------------------------------------------------

\needspace{9\baselineskip}\item
\textbf{ما التعريف الصحيح لـ«الأمر» (\textenglish{Prompt}) كما ورد في الدرس؟} \hfill \pts{0,75}
\begin{itemize}[leftmargin=1.6cm, labelsep=0.3cm, itemsep=2pt, topsep=4pt]
    \item[\emptybox\ \textbf{أ)}] قائمة الإعدادات التي نغيّرها في حساب الأداة.
    \item[\emptybox\ \textbf{ب)}] الإجابة التي تعرضها الأداة بعد معالجة الطلب.
    \item[\emptybox\ \textbf{ج)}] البرنامج الذي يشغّل الأداة داخل الحاسوب.
    \item[\correctbox\ \textbf{د)}] \correct{النص الذي نكتبه للأداة لنطلب منها شيئًا.}
\end{itemize}
\explanation{الأمر هو النص الذي نكتبه للأداة لنطلب منها شيئًا، وجودته تحدد إلى حد كبير جودة الإجابة.}

%------------------------------------------------

\needspace{9\baselineskip}\item
\textbf{سمع طالب عبارة «هندسة الأوامر» لأول مرة، وظنّ أنها تخص المهندسين فقط. ما المقصود بها فعلًا؟} \hfill \pts{0,75}
\begin{itemize}[leftmargin=1.6cm, labelsep=0.3cm, itemsep=2pt, topsep=4pt]
    \item[\emptybox\ \textbf{أ)}] برمجة أدوات الذكاء الاصطناعي بلغات برمجة معقدة، ولا يقوم بها إلا مهندسو الشركات التي تطوّر هذه الأدوات.
    \item[\emptybox\ \textbf{ب)}] حفظ قائمة من الأوامر الجاهزة وتكرارها كما هي في كل موقف، دون تكييفها مع المهمة أو مع المستوى المطلوب.
    \item[\correctbox\ \textbf{ج)}] \correct{مهارة كتابة أوامر واضحة ومنظمة للحصول على أفضل إجابة ممكنة من الأداة، ويمكن لأي مستخدم أن يتعلّمها.}
    \item[\emptybox\ \textbf{د)}] تخصص هندسي يدرس بناء الحواسيب التي تعمل عليها أدوات الذكاء الاصطناعي، ولا علاقة له بكتابة الطلبات.
\end{itemize}
\explanation{هندسة الأوامر مهارة كتابة أوامر واضحة ومنظمة للحصول على أفضل إجابة. لا تتطلب برمجة، بل تفكيرًا في ما تحتاجه الأداة لتفهم الطلب.}

%------------------------------------------------

\needspace{9\baselineskip}\item
\textbf{بدأت أستاذة أمرها بعبارة: «تصرّف كأستاذ لغة عربية خبير في تعليم المتوسط». أي عنصر من عناصر الأمر الجيد تمثّل هذه العبارة؟} \hfill \pts{0,75}
\begin{itemize}[leftmargin=1.6cm, labelsep=0.3cm, itemsep=2pt, topsep=4pt]
    \item[\emptybox\ \textbf{أ)}] المهمة، لأنها تحدد ما الذي نريد من الأداة أن تفعله بالضبط في الإجابة.
    \item[\correctbox\ \textbf{ب)}] \correct{الدور، لأنها تحدد من تكون الأداة وأي خبرة تعتمد عليها في الإجابة.}
    \item[\emptybox\ \textbf{ج)}] القيود، لأنها تحدد الشروط التي يجب احترامها مثل عدد الكلمات والمستوى.
    \item[\emptybox\ \textbf{د)}] التنسيق، لأنها تحدد الشكل الذي نريد أن تظهر به الإجابة على الشاشة.
\end{itemize}
\explanation{الدور يجيب عن سؤال: من تكون الأداة؟ وأي خبرة تعتمد؟ مثل: «تصرّف كأستاذ لغة عربية».}

%------------------------------------------------

\needspace{9\baselineskip}\item
\textbf{في أمر طويل، كتب أستاذ عبارة: «تلاميذي في السنة الثانية متوسط، ويجدون صعوبة في التمييز بين الحال والتمييز». ما العنصر الذي تمثّله هذه العبارة؟} \hfill \pts{0,75}
\begin{itemize}[leftmargin=1.6cm, labelsep=0.3cm, itemsep=2pt, topsep=4pt]
    \item[\emptybox\ \textbf{أ)}] التنسيق، لأنها تحدد شكل الإجابة وعدد أقسامها وطريقة ترتيب أجزائها.
    \item[\emptybox\ \textbf{ب)}] الدور، لأنها تحدد الخبرة التي يجب أن تتقمّصها الأداة أثناء كتابة الإجابة.
    \item[\emptybox\ \textbf{ج)}] المهمة، لأنها تحدد الفعل الأساسي الذي نطلبه من الأداة بشكل مباشر.
    \item[\correctbox\ \textbf{د)}] \correct{السياق، لأنها تقدّم معلومات تساعد الأداة على فهم الوضعية وتكييف إجابتها معها.}
\end{itemize}
\explanation{السياق يجيب عن سؤال: ما المعلومات التي تساعد الأداة على الفهم؟ مستوى التلاميذ وصعوباتهم معلومات سياقية.}

%------------------------------------------------

\needspace{9\baselineskip}\item
\textbf{أضافت أستاذة إلى أمرها: «في 150 كلمة، بلغة بسيطة، ودون استعمال مصطلحات أجنبية». إلى أي عنصر تنتمي هذه الإضافة؟} \hfill \pts{0,75}
\begin{itemize}[leftmargin=1.6cm, labelsep=0.3cm, itemsep=2pt, topsep=4pt]
    \item[\correctbox\ \textbf{أ)}] \correct{القيود، لأنها شروط يجب أن تحترمها الأداة في الإجابة.}
    \item[\emptybox\ \textbf{ب)}] الدور، لأنها تحدد الشخصية التي ستتقمّصها الأداة.
    \item[\emptybox\ \textbf{ج)}] السياق، لأنها تصف التلاميذ الذين ستُقدَّم لهم الإجابة.
    \item[\emptybox\ \textbf{د)}] المهمة، لأنها تحدد الدرس الذي ستشرحه الأداة بالضبط.
\end{itemize}
\explanation{القيود تجيب عن سؤال: ما الشروط التي يجب احترامها؟ مثل عدد الكلمات وبساطة اللغة وما يجب تجنّبه.}

%------------------------------------------------

\needspace{9\baselineskip}\item
\textbf{ختم أستاذ أمره بعبارة: «نظّم الإجابة في ثلاثة أجزاء: تعريف، ثم أمثلة، ثم خلاصة». أي عنصر يمثّله هذا الجزء؟} \hfill \pts{0,75}
\begin{itemize}[leftmargin=1.6cm, labelsep=0.3cm, itemsep=2pt, topsep=4pt]
    \item[\emptybox\ \textbf{أ)}] السياق، لأنه يشرح مستوى التلاميذ الذين سيقرؤون الشرح.
    \item[\emptybox\ \textbf{ب)}] الدور، لأنه يحدد الخبرة التي تعتمد عليها الأداة.
    \item[\correctbox\ \textbf{ج)}] \correct{التنسيق، لأنه يحدد الشكل الذي نريد أن تأتي عليه الإجابة.}
    \item[\emptybox\ \textbf{د)}] القيود، لأنه يحدد عدد الكلمات المسموح به في كل جزء.
\end{itemize}
\explanation{التنسيق يجيب عن سؤال: كيف نريد شكل الإجابة؟ مثل: تعريف ثم أمثلة ثم خلاصة، أو جدول، أو قائمة.}

%------------------------------------------------

\needspace{9\baselineskip}\item
\textbf{في الأمر: «أنا أستاذ متربص. تلاميذي في السنة الثانية متوسط. اشرح درس الحال»، ما الجزء الذي يمثّل المهمة؟} \hfill \pts{0,75}
\begin{itemize}[leftmargin=1.6cm, labelsep=0.3cm, itemsep=2pt, topsep=4pt]
    \item[\correctbox\ \textbf{أ)}] \correct{«اشرح درس الحال»، لأنه يحدد ما نريد أن تفعله الأداة.}
    \item[\emptybox\ \textbf{ب)}] لا توجد مهمة في هذا الأمر، لأنه لا يحدد عدد الكلمات.
    \item[\emptybox\ \textbf{ج)}] «أنا أستاذ متربص»، لأنه يحدد من يكتب الأمر للأداة.
    \item[\emptybox\ \textbf{د)}] «تلاميذي في السنة الثانية متوسط»، لأنه يحدد المستوى.
\end{itemize}
\explanation{المهمة تجيب عن سؤال: ماذا نريد أن تفعل الأداة؟ وهي هنا «اشرح درس الحال». أما المعلومات عن الأستاذ والتلاميذ فهي سياق.}

%------------------------------------------------

\needspace{9\baselineskip}\item
\textbf{حلّل طالب الأمر الآتي: «تصرّف كأستاذ جامعي في البلاغة. اشرح مفهوم التشبيه لطالب في السنة الأولى، في أقل من 200 كلمة، مع جدول يبيّن أركان التشبيه». أي توزيع للعناصر صحيح؟} \hfill \pts{0,75}
\begin{itemize}[leftmargin=1.6cm, labelsep=0.3cm, itemsep=2pt, topsep=4pt]
    \item[\correctbox\ \textbf{أ)}] \correct{«لطالب في السنة الأولى» سياق، و«في أقل من 200 كلمة» قيد، و«مع جدول يبيّن أركان التشبيه» تنسيق.}
    \item[\emptybox\ \textbf{ب)}] «لطالب في السنة الأولى» قيد، و«في أقل من 200 كلمة» تنسيق، و«مع جدول يبيّن أركان التشبيه» مهمة.
    \item[\emptybox\ \textbf{ج)}] «لطالب في السنة الأولى» مهمة، و«في أقل من 200 كلمة» دور، و«مع جدول يبيّن أركان التشبيه» سياق.
    \item[\emptybox\ \textbf{د)}] «لطالب في السنة الأولى» دور، و«في أقل من 200 كلمة» سياق، و«مع جدول يبيّن أركان التشبيه» قيد.
\end{itemize}
\explanation{في هذا المثال من الدرس: الدور «أستاذ جامعي في البلاغة»، والمهمة «اشرح مفهوم التشبيه»، والسياق «طالب في السنة الأولى»، والقيد «أقل من 200 كلمة»، والتنسيق «جدول».}

%------------------------------------------------

\needspace{9\baselineskip}\item
\textbf{كتبت أستاذة: «تصرّف كمدقق لغوي. صحّح الأخطاء في النص الآتي، فهو لتلميذ في السنة الرابعة متوسط. اعرض التصحيح في جدول من ثلاثة أعمدة: الخطأ، الصواب، القاعدة». أي عنصر من العناصر الخمسة غائب عن هذا الأمر؟} \hfill \pts{0,75}
\begin{itemize}[leftmargin=1.6cm, labelsep=0.3cm, itemsep=2pt, topsep=4pt]
    \item[\emptybox\ \textbf{أ)}] الدور: لأن عبارة «تصرّف كمدقق لغوي» لا تكفي، ويجب ذكر اسم الأداة المستعملة في الأمر.
    \item[\correctbox\ \textbf{ب)}] \correct{القيود: لم تحدد شروطًا مثل عدم تغيير أسلوب التلميذ أو الاكتفاء بالأخطاء النحوية والإملائية.}
    \item[\emptybox\ \textbf{ج)}] السياق: لأن مستوى التلميذ لا يُعدّ معلومة سياقية، بل يُعدّ تنسيقًا للإجابة النهائية.
    \item[\emptybox\ \textbf{د)}] التنسيق: لأن الجدول ليس شكلًا للإجابة، بل هو مهمة إضافية يجب كتابتها في أمر مستقل.
\end{itemize}
\explanation{الأمر يحتوي على الدور (مدقق لغوي) والمهمة (صحّح) والسياق (الرابعة متوسط) والتنسيق (جدول). الغائب هو القيود: الشروط التي يجب احترامها.}

%------------------------------------------------

\needspace{9\baselineskip}\item
\textbf{اقترح الدرس بناء الأمر خطوة بخطوة انطلاقًا من أمر عام. أي ترتيب للدرجات يوافق ما رأيناه؟} \hfill \pts{0,75}
\begin{itemize}[leftmargin=1.6cm, labelsep=0.3cm, itemsep=2pt, topsep=4pt]
    \item[\emptybox\ \textbf{أ)}] إضافة التنسيق $\leftarrow$ أمر عام $\leftarrow$ إضافة السياق $\leftarrow$ إضافة المهمة والقيود.
    \item[\correctbox\ \textbf{ب)}] \correct{أمر عام $\leftarrow$ إضافة السياق $\leftarrow$ إضافة المهمة والقيود $\leftarrow$ إضافة التنسيق.}
    \item[\emptybox\ \textbf{ج)}] إضافة القيود $\leftarrow$ إضافة التنسيق $\leftarrow$ أمر عام $\leftarrow$ إضافة السياق والدور.
    \item[\emptybox\ \textbf{د)}] أمر عام $\leftarrow$ إضافة التنسيق $\leftarrow$ حذف السياق $\leftarrow$ إضافة المهمة فقط.
\end{itemize}
\explanation{في الدرس: الدرجة 1 أمر عام («اشرح الحال»)، ثم نضيف السياق، ثم المهمة والقيود، ثم التنسيق، فتتحسن جودة الإجابة في كل درجة.}

%------------------------------------------------

\needspace{9\baselineskip}\item
\textbf{كتب طالب: «اشرح البلاغة». ما الخطأ الشائع الذي وقع فيه حسب الدرس؟} \hfill \pts{0,75}
\begin{itemize}[leftmargin=1.6cm, labelsep=0.3cm, itemsep=2pt, topsep=4pt]
    \item[\emptybox\ \textbf{أ)}] أمر يخلط عدة مهام مختلفة، فيجب تقسيمه إلى أوامر متتابعة.
    \item[\emptybox\ \textbf{ب)}] لا خطأ فيه، لأن الأداة تعرف وحدها مستوى الطالب وهدفه.
    \item[\correctbox\ \textbf{ج)}] \correct{أمر قصير جدًا: لا يحدد المستوى ولا الهدف، فتأتي الإجابة عامة.}
    \item[\emptybox\ \textbf{د)}] أمر طويل جدًا: يحتوي على تفاصيل كثيرة تُربك الأداة وتبطئها.
\end{itemize}
\explanation{«اشرح البلاغة» مثال الدرس على الأمر القصير جدًا: لا يحدد المستوى ولا الهدف، فلا تعرف الأداة ماذا يريد الطالب بالضبط.}

%------------------------------------------------

\needspace{9\baselineskip}\item
\textbf{طلب تلميذ من الأداة: «اكتب شيئًا جميلًا». ما المشكلة الأساسية في هذا الأمر؟} \hfill \pts{0,75}
\begin{itemize}[leftmargin=1.6cm, labelsep=0.3cm, itemsep=2pt, topsep=4pt]
    \item[\emptybox\ \textbf{أ)}] أمر يخلط بين الدور والتنسيق، فيجب الفصل بينهما في جملتين.
    \item[\emptybox\ \textbf{ب)}] أمر مكتوب بلغة عربية، والأدوات لا تفهم إلا الأوامر الإنجليزية.
    \item[\emptybox\ \textbf{ج)}] أمر يحتوي على قيود كثيرة جدًا تمنع الأداة من الإبداع بحرية.
    \item[\correctbox\ \textbf{د)}] \correct{أمر غامض: لا يحدد ما المطلوب، لا نوع النص ولا موضوعه ولا طوله.}
\end{itemize}
\explanation{«اكتب شيئًا جميلًا» مثال الدرس على الأمر الغامض: كلمة «شيء» لا تحدد نوع المحتوى ولا موضوعه ولا جمهوره.}

%------------------------------------------------

\needspace{9\baselineskip}\item
\textbf{كتبت أستاذة في أمر واحد: «اشرح درس الحال، واكتب عشرين تمرينًا، وصحّحها، وترجم الشرح إلى الفرنسية، وأعدّ اختبارًا نهائيًا مع سلّم تنقيط». ما النصيحة الأنسب لها حسب الدرس؟} \hfill \pts{0,75}
\begin{itemize}[leftmargin=1.6cm, labelsep=0.3cm, itemsep=2pt, topsep=4pt]
    \item[\emptybox\ \textbf{أ)}] حذف كل المهام والاكتفاء بطلب الاختبار النهائي، لأن الأداة لا تستطيع أبدًا القيام بأكثر من مهمة واحدة في المحادثة.
    \item[\correctbox\ \textbf{ب)}] \correct{تقسيم هذه المهمة الكبيرة إلى أوامر متتابعة، لكل أمر مهمة واحدة واضحة، ثم مراجعة كل إجابة قبل الانتقال إلى التالية.}
    \item[\emptybox\ \textbf{ج)}] إضافة مهام أخرى إلى الأمر نفسه، لأن الأداة تفهم الأوامر الطويلة التي تحتوي على مهام كثيرة أفضل من الأوامر القصيرة.
    \item[\emptybox\ \textbf{د)}] الإبقاء على الأمر كما هو، لأن جمع المهام كلها في أمر واحد يوفّر الوقت ويعطي دائمًا نتيجة أدق من الأوامر المنفصلة.
\end{itemize}
\explanation{من الأخطاء الشائعة في الدرس جمع عدة مهام مختلفة في أمر واحد. الأفضل تقسيم المهام الكبيرة إلى أوامر متتابعة.}

%------------------------------------------------

\needspace{9\baselineskip}\item
\textbf{لم تكن أستاذة راضية عن أول إجابة قدّمتها الأداة لشرح درس النعت. ما التصرف الذي يوافق الدرس؟} \hfill \pts{0,75}
\begin{itemize}[leftmargin=1.6cm, labelsep=0.3cm, itemsep=2pt, topsep=4pt]
    \item[\emptybox\ \textbf{أ)}] إعادة إرسال الأمر نفسه حرفيًا عدة مرات دون أي تعديل، لأن تكرار الأمر هو الطريقة الوحيدة لتحسين الإجابة.
    \item[\emptybox\ \textbf{ب)}] التخلي نهائيًا عن استعمال الأداة، لأن الإجابة غير المرضية دليل على أن الأداة لا تصلح لدروس اللغة العربية.
    \item[\correctbox\ \textbf{ج)}] \correct{تعديل الأمر بإضافة ما ينقصه، مثل السياق أو القيود أو التنسيق، ثم إعادة المحاولة ومقارنة الإجابة الجديدة بالسابقة.}
    \item[\emptybox\ \textbf{د)}] قبول الإجابة الأولى كما هي، لأن الأداة تقدّم دائمًا أفضل ما لديها من أول محاولة، ولا تتغير إجابتها بتغيير الأمر.
\end{itemize}
\explanation{«قبول أول إجابة» من الأخطاء الشائعة. يمكن دائمًا تعديل الأمر وإعادة المحاولة، فالأمر يُبنى تدريجيًا ويُحسَّن بعد قراءة الإجابة.}

%------------------------------------------------

\needspace{9\baselineskip}\item
\textbf{قال الدرس: «الأداة لا تعرف ما في ذهنك». ما النتيجة العملية لهذه الفكرة عند كتابة الأمر؟} \hfill \pts{0,75}
\begin{itemize}[leftmargin=1.6cm, labelsep=0.3cm, itemsep=2pt, topsep=4pt]
    \item[\emptybox\ \textbf{أ)}] لا داعي لتحديد المستوى، لأن الأداة تكتشف مستوى المستخدم من طريقة كتابته.
    \item[\emptybox\ \textbf{ب)}] يجب أن نكتب الأمر بالإنجليزية دائمًا، لأن الأداة لا تفهم ما في أذهان العرب.
    \item[\emptybox\ \textbf{ج)}] يكفي أن نكتب كلمة واحدة، لأن الأداة تخمّن الباقي من سجلّ استعمالنا السابق.
    \item[\correctbox\ \textbf{د)}] \correct{يجب أن يكون الأمر واضحًا وكاملًا، لأن الأداة لا تعرف إلا ما نكتبه لها فعلًا.}
\end{itemize}
\explanation{بما أن الأداة لا تعرف إلا ما نكتبه لها، فكل معلومة مهمة (المستوى، الهدف، الشكل...) يجب أن تُذكر صراحة في الأمر.}

%------------------------------------------------

\needspace{9\baselineskip}\item
\textbf{أرادت أستاذة أمثلة على المفعول لأجله لتلاميذها. أي أمر هو الأفضل حسب عناصر الأمر الجيد؟} \hfill \pts{0,75}
\begin{itemize}[leftmargin=1.6cm, labelsep=0.3cm, itemsep=2pt, topsep=4pt]
    \item[\correctbox\ \textbf{أ)}] \correct{تصرّف كأستاذ لغة عربية. اقترح عشر جمل فيها مفعول لأجله لتلاميذ السنة الثالثة متوسط، من الحياة اليومية، وضع خطًا تحت المفعول لأجله في كل جملة.}
    \item[\emptybox\ \textbf{ب)}] اقترح جملًا فيها مفعول لأجله، واكتب أي عدد تراه مناسبًا من الجمل بأي مستوى، فالمهم أن تكون الجمل كثيرة وطويلة قدر المستطاع.
    \item[\emptybox\ \textbf{ج)}] المفعول لأجله. جمل. أمثلة. نحو. متوسط. لغة عربية. تلاميذ. أستاذ. درس. قاعدة. تمارين. إعراب. شرح. كثير.
    \item[\emptybox\ \textbf{د)}] أنت تعرف ما أريد، فأعطني ما يحتاجه تلاميذي في درس الغد دون أن أشرح لك التفاصيل، فأنت أدرى بها مني.
\end{itemize}
\explanation{الأمر الأول يجمع الدور (أستاذ لغة عربية) والمهمة (اقترح جملًا) والسياق (الثالثة متوسط) والقيود (عشر جمل من الحياة اليومية) والتنسيق (خط تحت المفعول لأجله).}

%------------------------------------------------

\needspace{9\baselineskip}\item
\textbf{تريد أستاذة أن تحصل على الإجابة في جدول يقارن بين أنواع الخبر. أي عنصر عليها أن تضيفه إلى أمرها؟} \hfill \pts{0,75}
\begin{itemize}[leftmargin=1.6cm, labelsep=0.3cm, itemsep=2pt, topsep=4pt]
    \item[\emptybox\ \textbf{أ)}] الدور: مثل «تصرّف كأستاذ نحو»، لأنه يحدد شكل الجدول بدقة.
    \item[\correctbox\ \textbf{ب)}] \correct{التنسيق: مثل «اعرض الإجابة في جدول من عمودين: نوع الخبر، ومثال».}
    \item[\emptybox\ \textbf{ج)}] السياق: مثل «تلاميذي في المتوسط»، لأنه يحدد عدد أعمدة الجدول.
    \item[\emptybox\ \textbf{د)}] المهمة: مثل «اشرح الخبر»، لأنها تكفي لتظهر الإجابة في جدول.
\end{itemize}
\explanation{شكل الإجابة (جدول، قائمة، فقرات...) يُحدَّد في عنصر التنسيق.}

%------------------------------------------------

\needspace{9\baselineskip}\item
\textbf{طُلب من طالب أن يحدد أي العبارات الآتية يوافق خلاصة درس هندسة الأوامر. أي عبارة صحيحة؟} \hfill \pts{0,75}
\begin{itemize}[leftmargin=1.6cm, labelsep=0.3cm, itemsep=2pt, topsep=4pt]
    \item[\emptybox\ \textbf{أ)}] الأمر الجيد هو الأطول دائمًا، ولا يحتاج إلى أي تعديل بعد كتابته للمرة الأولى.
    \item[\emptybox\ \textbf{ب)}] الأمر الجيد يُكتب مرة واحدة ويُستعمل كما هو مع كل الأدوات وكل المهام والمستويات.
    \item[\correctbox\ \textbf{ج)}] \correct{الأمر الجيد يعطي الأداة كل ما تحتاجه، ونبنيه تدريجيًا ونحسّنه بعد قراءة الإجابة.}
    \item[\emptybox\ \textbf{د)}] الأمر الجيد يكتفي بالمهمة وحدها، لأن الدور والسياق والقيود والتنسيق كلها حشو.
\end{itemize}
\explanation{خلاصة الدرس: الأمر الجيد يعطي الأداة كل ما تحتاجه، عناصره خمسة، ونبنيه تدريجيًا ونحسّنه بعد قراءة الإجابة.}

%------------------------------------------------

\needspace{9\baselineskip}\item
\textbf{تريد أستاذة تحويل أبيات شعرية إلى نثر لتلاميذها. أي أمر يوافق نموذج «الأوامر المفيدة لأستاذ اللغة العربية» في الدرس؟} \hfill \pts{0,75}
\begin{itemize}[leftmargin=1.6cm, labelsep=0.3cm, itemsep=2pt, topsep=4pt]
    \item[\emptybox\ \textbf{أ)}] اكتب لي شيئًا عن هذه الأبيات، وأضف ما تشاء من المعلومات حتى لو لم تكن لها علاقة مباشرة بالنص أو بمستوى التلاميذ.
    \item[\emptybox\ \textbf{ب)}] أبيات، نثر، تلاميذ. افعل ما تراه مناسبًا وغيّر المعنى إن لزم الأمر، فالمهم أن يكون النص النهائي طويلًا ومنمّقًا.
    \item[\emptybox\ \textbf{ج)}] ترجم هذه الأبيات إلى الإنجليزية ثم أعدها إلى العربية عدة مرات، فهذه أفضل طريقة لتحويل أي شعر إلى نثر بسيط.
    \item[\correctbox\ \textbf{د)}] \correct{حوّل هذه الأبيات إلى نثر بسيط يفهمه تلاميذ السنة الأولى متوسط، دون تغيير المعنى، ثم اذكر الفكرة العامة في سطر واحد.}
\end{itemize}
\explanation{الأمر الأول يحدد المهمة (حوّل إلى نثر) والسياق (الأولى متوسط) والقيد (دون تغيير المعنى) وشكل الإضافة المطلوبة، وهو قريب من أمثلة الدرس.}

\end{enumerate}

\vfill
\begin{center}
{\khatt بالتوفيق}
\end{center}

\end{document}
